(1) 薄膜局部受力为张力乘横向Laplacian，故 \(\sigma w_{tt}=T[w_{rr}+r^{-1}w_r+r^{-2}w_{\theta\theta}]\)，波速 \(c=\sqrt{T/\sigma}\)。设 \(w=R(r)e^{im\theta}e^{-i\omega t}\)，有 \(R''+R'/r+(k^2-m^2/r^2)R=0\)，\(k=\omega/c\)。

(2) 对非零 \(k\)，\(R=A J_m(kr)+B Y_m(kr)\)。施加 \(R(a)=R(b)=0\) 的线性齐次方程存在非零解当且仅当 \(\boxed{J_m(ka)Y_m(kb)-Y_m(ka)J_m(kb)=0}\)。相应本征频率 \(\omega_{m,n}=c k_{m,n}\)，\(k_{m,n}>0\) 为该行列式依次正零点。

(3) \(m=0\) 模式角向与 \(\theta\) 无关，角简并为1；\(m\ge1\) 的 \(\cos m\theta,\sin m\theta\) 独立且等频，简并为2。两个圆边均施加Dirichlet零位移，均匀平移不满足边界；由能量积分 \(\int|\nabla w|^2dA>0\) 对非零膜位移成立，故没有零频非平凡模。

(4) 自伴径向算子为 \(-r^{-1}(rR')'+m^2R/r^2\)。对不同径向本征值 \(k_n^2\ne k_p^2\)，将两本征方程交叉相乘后积分并分部，边界项因 \(R(a)=R(b)=0\) 消失，故 \(\int_a^b rR_nR_pdr=0\)。全二维内积还含 \(\int_0^{2\pi}d\theta\)。薄环固定 \(m\) 且 \(\ell=b-a\ll a\) 时 \(k_{m,n}\sim n\pi/\ell\)，弯曲边界与角向项给次阶修正。
